\chapter{Discussion}
\label{ch:discussion}

\section{Interpretation of Findings by Research Question}
\label{sec:discussion-rq}

\subsection{RQ1: Overlap with the Baseline}

Across all 300 profiles and three strategies, LLM-generated recommendations showed meaningful but partial agreement with the content-based cosine-similarity baseline (mean overlap ranging from 0.517 to 0.687 depending on strategy, against a maximum possible overlap of 1.0). This confirms that an 8-billion-parameter, free-tier LLM can approximate the output of a purpose-built content-based recommender to a substantial degree without any task-specific training, supporting the premise in Section~1.1 that LLMs are a viable alternative recommendation mechanism. It is not, however, evidence that the LLM and the baseline are interchangeable: an overlap of 0.687 at its highest (zero-shot) still means, on average, roughly one of the baseline's three reference cards is absent from the LLM's top three. Because the baseline itself is a fixed-formula approximation of ``good fit'' rather than ground truth in an absolute sense, a lower overlap score is not necessarily evidence of a worse recommendation --- it may equally reflect the LLM incorporating profile information (a stated goal, a qualitative preference) that the baseline's eight-dimensional feature vector cannot represent. This is an inherent limitation of using a deterministic baseline as the sole reference point, discussed further in Section~\ref{sec:limitations-discussion}.

\subsection{RQ2: Effect of Prompting Strategy}

The result that zero-shot prompting significantly outperformed both structured and few-shot prompting on recommendation-baseline overlap and top-1 accuracy (Section~4.2) runs counter to the general expectation, drawn from the wider prompt-engineering literature (Section~2.2), that providing worked examples (few-shot) or an explicit schema (structured) should improve task performance. Three plausible, non-exclusive explanations are considered here.

First, both the structured and few-shot prompts are substantially longer than the zero-shot prompt --- the structured prompt adds a full JSON schema specification, and the few-shot prompt adds two complete worked examples on top of that same instruction set. At 8 billion parameters, Llama 3.1 8B has comparatively limited capacity to maintain attention over long, multi-part instructions relative to larger frontier models; it is plausible that the extra instructional content in these two conditions displaces some of the model's effective attention on the specific profile it is meant to be reasoning about, a phenomenon informally consistent with the ``lost in the middle'' effect reported for long-context prompting in other studies, though this study did not directly test that mechanism and it should be treated as a plausible interpretation rather than a demonstrated cause.

Second, the two worked examples in the few-shot prompt necessarily represent only two of the eight archetypes in this study's dataset (a student-type profile and a frequent-traveller-type profile). The archetype-level breakdown in Table~4.2 shows the few-shot strategy performing at or near zero top-1 accuracy for several archetypes distant from the two worked examples (e.g.\ \texttt{family\_grocery\_heavy}, \texttt{balanced\_cashback}), which is consistent with the model anchoring more heavily on the specific examples shown than on the general instruction, rather than generalising the underlying task.

Third, the factual-correctness results (Section~4.4) show few-shot prompting generating the largest volume of specific numeric claims while being the least accurate on those claims. This suggests the worked examples may be teaching the model a superficial pattern --- ``state a specific fee or reward-rate figure to sound authoritative'' --- rather than the underlying reasoning process, a documented risk of few-shot prompting when the examples are not sufficiently representative of the full input distribution.

These findings should not be read as a general claim that zero-shot prompting is superior to few-shot or structured prompting; they are specific to this model scale, this task, and this particular set of worked examples, and are precisely the kind of task-dependent result the mixed literature in Section~2.2 would predict.

\subsection{RQ3: Factual Reliability and Quality of Stated Reasoning}

The factual-correctness results (Section~4.4) show that requiring an explicit output schema (structured prompting) produced the most factually reliable responses (73.7\% entirely error-free), even though structured prompting did not produce the best recommendation-quality scores. This is the clearest demonstration in this study's results of why a single-metric evaluation would be inadequate for this task: a marker evaluating only recommendation overlap would conclude zero-shot is the best strategy outright, while a marker evaluating only factual reliability would reach the opposite conclusion for structured prompting. The four-criteria framework adopted in Section~3.4.3, motivated directly by the gap identified in the literature review (Section~2.6), is what makes this trade-off visible.

The explanation-quality results (Section~\ref{sec:explanation-consistency}) sharpen this picture further, and produce the single most important finding of this dissertation: \textbf{the strategy judged to write the best-quality explanations (few-shot, mean 2.907/3) is simultaneously the least factually reliable (24.0\% error-free) and the least accurate against the baseline (mean overlap 0.539)}. A judge --- whether an LLM, as here, or plausibly a human reader without access to ground truth --- rating explanation quality from the text alone would rank few-shot's output as the best of the three strategies. Only by separately, deterministically checking the specific factual claims against the real card catalogue does it become clear that this apparent quality is not matched by accuracy. This is a direct, empirical illustration of a risk that is easy to state abstractly but hard to demonstrate concretely: a fluent, well-structured, confident-sounding explanation is not evidence that an LLM's recommendation is correct, and evaluating explanation quality in isolation --- without an independent factual check --- would have reached the wrong conclusion about which strategy to prefer.

\subsection{RQ4: Consistency}

The consistency results (Section~\ref{sec:consistency-results}) show zero-shot as numerically the most stable strategy (mean Jaccard 0.826) and structured as the least stable (0.776), with few-shot in between (0.808) --- though, unlike the RQ2 and RQ3 results, this difference is not statistically significant on the 24-profile stratified sample used ($p = 0.68$), so it should be read as a trend rather than a confirmed effect. Taken together with RQ2 and RQ3, an interesting pattern emerges rather than one strategy simply being best: zero-shot is both the most accurate against the baseline \textit{and} numerically the most consistent, which is internally coherent (a model reasoning more directly from the profile, without extra instructional content to weigh, might plausibly be both more accurate and more repeatable). Structured prompting, by contrast, is the most factually reliable but numerically the least consistent of the three --- a response that gets the facts right more often is not necessarily the response that stays the same each time you ask. This is a genuine, if statistically tentative, three-way tension across the criteria rather than a single dominant strategy, and reinforces this dissertation's central methodological argument (Section~2.6) that no single evaluation criterion tells the whole story for this kind of task.

\section{Relation to the Literature}

The finding that prompting strategy has a large, statistically significant effect on this task (Section~4.2) is consistent with the general premise of the prompt-engineering literature (Schulhoff et al., 2024) that prompt design is a first-class variable rather than an implementation detail. However, the specific direction of the effect --- zero-shot outperforming few-shot --- diverges from the more common finding in general NLP benchmarks that few-shot prompting improves performance (Section~2.2), reinforcing this dissertation's positioning (Section~2.6) that recommendation tasks with a real, checkable product catalogue behave differently from the classification and QA benchmarks that dominate the existing zero-shot/few-shot comparison literature. The factual-correctness results also extend the hallucination concerns raised generally in the LLM-recommendation survey literature (Zhao et al., 2024; Wu et al., 2024) with a concrete, measured error rate specific to this domain, rather than a general acknowledgement that hallucination is possible.

The explanation-quality/factual-correctness divergence for few-shot prompting (Section~\ref{sec:discussion-rq}) also connects directly to a specific, named risk in the LLM-as-judge literature: Zheng et al.\ (2023) document \textit{verbosity bias}, where a judge (model or human) tends to rate longer, more detailed responses more favourably regardless of their actual correctness. The individual-level validity check in Section~\ref{sec:judge-validity} provides direct evidence consistent with exactly this risk: judge scores correlate weakly but significantly \textit{negatively} with individual responses being factually error-free ($r = -0.173$, $p < 0.0001$), meaning the judge model was not, in this study, a reliable proxy for factual accuracy even setting aside the strategy-level pattern --- few-shot's worked examples appear to have taught the model to write more detailed, specific-sounding justifications, and it is exactly this detail that the judge rates highly and that the factual checker subsequently shows to be frequently wrong. This is an argument for treating explanation-quality scores (from any judge, automated or human) as informative but insufficient on their own, and for pairing them with an independent, deterministic factual check wherever the underlying facts are checkable, as this study did.

\section{Limitations}
\label{sec:limitations-discussion}

\textbf{Synthetic data.} All 300 user profiles are synthetic, calibrated against aggregate ONS expenditure statistics rather than drawn from real individual spending records (Section~3.2.2), because no such records are legally available for this kind of research (UK GDPR, Financial Services and Markets Act 2000). ONS calibration anchors the plausibility of the aggregate distributions used but cannot capture the irregular, bursty, or life-event-driven spending patterns a real consumer would exhibit, so findings about which prompting strategy performs best may not transfer perfectly to real users with messier spending histories.

\textbf{Single model.} All recommendations were generated by one model (Llama 3.1, 8B parameters). The explanation for the zero-shot finding proposed in Section~\ref{sec:discussion-rq} --- that limited attention capacity at this model scale disadvantages longer prompts --- is a testable hypothesis, not a confirmed result, because this study did not repeat the experiment with a larger model to check whether the effect persists, diminishes, or reverses at greater scale.

\textbf{Baseline as reference point, not ground truth.} As discussed in Section~\ref{sec:discussion-rq}, the cosine-similarity baseline (Section~3.3) is itself a simplified, eight-dimensional approximation of ``good fit,'' not an independently validated gold standard. Overlap with the baseline measures agreement with one particular reasonable method, not objective recommendation quality in an absolute sense.

\textbf{LLM-as-judge for explanation quality.} Using a second LLM to score explanation quality (Section~3.4.3) is grounded in the literature (Zheng et al., 2023) and specifically designed to reduce self-enhancement bias by using a distinct, larger judge model, but it still substitutes model judgement for the human judgement originally specified in the project's methodology, and inherits other documented LLM-judge biases (position and verbosity bias) that a human rater would not share.

\textbf{Deterministic factual-correctness scope.} The factual-correctness checker (Section~3.4.3, \texttt{src/factual\_correctness.py}) can only verify claims that are specific, numeric, or exactly named against the card catalogue; it cannot assess vague or qualitative claims (e.g.\ ``this card offers good value''). The reported error rates (Section~4.4) are therefore a lower bound on the true rate of factual inaccuracy in the responses, not a complete audit of every claim made.

\textbf{Consistency sampling scope.} The consistency evaluation (Section~3.4.4) uses a stratified 24-profile subsample rather than all 300 profiles, for the practical rate-limit reasons given there; findings about stability generalise to the sampled archetypes but were not verified across every individual profile.

\textbf{Infrastructure constraints on methodology.} The free Groq API tier's rate and daily-request limits (Section~3.4.1) meant that both the primary 900-call experiment and the subsequent zero-shot/few-shot correction rerun (Section~3.4.2) each took multiple days of wall-clock time to complete, constraining how many additional experimental conditions (e.g.\ additional models, additional prompting strategies, a full-N consistency check) could practically be included within the scope of this dissertation.
